% page 294 of the facsimile (image p-293 of the scan)
% block 154.9 x 237.7 mm ; left margin 8.0 mm
\pgc{-12.700mm}{0.000mm}{
\l{286}
\l{}
\l{}
\l{\sou{kolodiono}. - Solvuro ek pulvero-kotono en etero, alkchol-}
\l{\cel{3}izita, uzata kirurgiale kom aglutinivo, e fotografale}
\l{\cel{3}por igar adheriva, sur la vitro-plaki, la strato de}
\l{\cel{3}iodido. - DEFIRS.}
\l{}
\l{\sou{kolofono}.\cel{2}Materio rezinoza, ye qua onu fricionas la krini}
\l{\cel{3}di la arketo por ke olu ne glitez sur la kordi di la}
\l{\cel{3}violino, di la violoncelo, e c. - DEFIRS.}
\l{}
\l{\sou{koloido}.\cel{2}(\sou{kemio}).Kompozajo ne-kristaligebla e ne-dializebla.}
\l{\cel{3}- DEF.}
\l{}
\l{\sou{kolokar}. (\sou{trans.})\cel{2}Prestar (mediate o nemediate) pekunio por}
\l{\cel{3}recevor de olu revenuo. - IS.}
\l{}
\l{\sou{kolombo}. (\sou{zool.}) Ucelo, meza inter la \textquotedbl{}galinacei\textquotedbl{} e la \textquotedbl{}pa-}
\l{\cel{3}seri\textquotedbl{}, di qua la pektoro-plumi esas tre koloroza. -}
\l{\cel{3}L.\cel{1}\sou{columba}.}
\l{}
\l{\sou{kolono}.-\cel{2}I. (\sou{arkitekt.})\cel{2}Pilastro ek quadro, ek marmoro,}
\l{\cel{3}qua finas per kapitelo, uzata kom subtenilo o kom orn-}
\l{\cel{3}ivo di edifico. - II. (\sou{anat.}) Parto di la grosa intes-}
\l{\cel{3}tino, qua sequas la ceko e su extensas til la rektumo.}
\l{\cel{3}- DEFIRS.}
\l{}
\l{\sou{kolonelo}. La viro qua, en armeo, komandas regimento. -DEFIS.}
\l{}
\l{\sou{kolonio}. I. Establisuro fondita da naciono di qua ula quanto}
\l{\cel{3}de habitanti iras aden regiono stranjera, ed ibe su ins-}
\l{\cel{3}talas fixe, sen cesar apartenar a la metropolo. - II. Ko-}
\l{\cel{3}lonio agrokultivala : establisuro rurala, destinita okupor}
\l{\cel{3}laboro povri, yuna karcerani, ed emendar la sulo. - III.}
\l{\cel{3}Grupo de individui ek un naciono, instalita fixe r lando}
\l{\cel{3}stranjera, en rezido-regiono komuna. - DEFIRS.}
\l{}
\l{\sou{koloquinto}. (\sou{bot.})\cel{2}Varietato di\cel{2}kukombro. - L.\cel{1}\sou{cucumis}}
\l{\cel{3}\sou{colycynthis}.- DEFIRS.}
\l{}
\l{\sou{kolor}o.I. Impreso partikulara, quan agas, sur la vid-organo,}
\l{\cel{3}ta o ca ek la elementi aden qui su deskompozas lumo-radio,}
\l{\cel{3}forsendata da surfaco qua reflektas ta elemento absor--}
\l{\cel{3}ante la cetera. - II. Proprajo di ula korpi pri tale im-}
\l{\cel{3}presar. - III. (\sou{metaf}.)Karaktero semblanta di la kozi.}
\l{\cel{3}- DEFIS.}
\l{}
\l{\sou{kolorito}. Efekto qua rezultas de la selekto, de la uzo di}
\l{\cel{3}ula farbi en pikturo.-\cel{2}DFIRS.}
\l{}
\l{\sou{koloso}.- I. Statuo di qua la dimensioni esas grandega.}
\l{\cel{3}II. Homo, animalo, di qua la dimensiono esas grandega.}
\l{\cel{3}III. (\sou{metaf.}) La homo qua esas grandega pri ca o ta ek}
\l{\cel{3}lua qualesi. - DEFIRS.}
}
